\documentclass[10pt,a4paper]{amsart}
\usepackage[margin=19mm]{geometry}
\usepackage{amsmath,amssymb,mathtools}
\usepackage[scr=rsfs,cal=euler]{mathalfa}
\usepackage{xfrac}
\usepackage[unicode,hidelinks,bookmarks=false]{hyperref}
\newcommand{\ie}{i.e.\ }
\newcommand{\cf}{cf.\ }
\newcommand{\resp}{respectively\ }
\newcommand{\Subsection}{\subsection}
\providecommand{\nobreakdash}{\nobreak\hbox{-}}
\setlength{\emergencystretch}{2em}
\multlinegap0pt
\usepackage{bm}
\usepackage{bbm}
\usepackage{dsfont}
\usepackage{xspace}
\usepackage{cancel}
\usepackage{mathtools}
\usepackage{amsthm}
\makeatletter
\@ifundefined{theorem}{\newtheorem{theorem}{Theorem}}{}
\@ifundefined{lem}{\newtheorem{lem}[theorem]{Lemma}}{}
\makeatother
\title[On generating sets of left-compressed intersecting families]{On generating sets of left-compressed intersecting families}
\author{Tuan Nguyen, Thi Nguyen, Thu Tran}
\date{Source: arXiv:2510.05666v3 (CC BY 4.0)}
\begin{document}
\maketitle

\noindent\emph{Source and licence.} On generating sets of left-compressed intersecting families. Version arXiv:2510.05666v3 (CC BY 4.0), distributed under the Creative Commons Attribution 4.0 International licence, as recorded in the arXiv licence metadata for this exact version and shown on the abstract page https://arxiv.org/abs/2510.05666v3. Full rights notice: licenses/arxiv-2510.05666-CC-BY-4.0.txt.\par\smallskip

\noindent\textit{Editorial scope.} Counted unit: the Lem whose source label is \texttt{lem2} in the article (source0.tex lines 548--552, source label \texttt{lem2}), delivered together with the article's own complete original proof of it (source0.tex lines 553--572, one \texttt{proof} environment of 20 lines). No mathematics is written, repaired or re-derived: statement and proof are the article's own text line for line. Required context retained uncounted: none. Every definition, display and label of those lines is retained unchanged. Omitted from this package and not redistributed: 1--547, 573--624. Nothing inside a retained excerpt is omitted or altered: every retained body is delivered exactly as the article's lines are written, so no omission receipt of any kind accompanies this package. Citations: the retained excerpts carry no citation command at all, so no bibliography entry of the article is transcribed and no reference is redistributed. Reference adaptation: none was needed and none was made; every reference inside the delivered unit resolves inside it, so no unresolved reference or citation is produced. Printed slips kept literal: no printed word, symbol, hypothesis or number of the counted statement or of its proof was altered. Layout and class: the article's own class is \texttt{amsart} with packages including etex, amsmath, amssymb, amsthm, bbm, mathtools, comment, enumitem, hyperref, euscript, tikz, babel, adjustbox, cmtiup; the wrapper is the lane's pinned \texttt{amsart} editorial wrapper at 10pt on A4, which restates the theorem-like environments the retained text needs and adds the article's own macro definitions that the retained text needs (none), with extra wrapper packages (bm, bbm, dsfont, xspace, cancel, mathtools). The delivered numbering is the unit's own. Assets: no retained span contains a figure environment, a graphics-inclusion command, a PDF-page inclusion, a table or a TikZ picture; no image file is copied into this package, the package declares no asset dependency and no image file is redistributed with it. Licence: version arXiv:2510.05666v3 of this article is distributed under the Creative Commons Attribution 4.0 International licence, as recorded in the arXiv licence metadata for this exact version and shown on the abstract page https://arxiv.org/abs/2510.05666v3. A full rights notice is delivered at licenses/arxiv-2510.05666-CC-BY-4.0.txt. Characters: every retained excerpt is pure ASCII, so the delivered engine, XeLaTeX with its default Latin Modern text font, reports no missing character.
\begin{lem}\label{lem2} Let $G,H \in 2^{[n]}$. Then, the following statements are equivalent:\\
(i) $G$ and $H$ are strongly intersecting;\\
(ii) $\mathcal{F}(G)$ and $\mathcal{F}(H)$ are cross-intersecting;\\
(iii) There exists $1 \leq l \leq n$ such that $\mu_{G}(l)+\mu_{H}(l) > l$.
\end{lem}

\begin{proof} $(i) \Rightarrow (ii)$\\
Take $A \in \mathcal{F}(G), B \in \mathcal{F}(H)$. Then, $A \preceq G$ and $B \preceq H$. Let $A'$ be the set of the first $|G|$ elements of $A$ and $B'$ be the set of the first $|H|$ elements of $B$. Then, $A' \leq G, B' \leq H$. Since $G \sim_{si} H$, we get $A' \cap B' \neq \emptyset$ and therefore $A \cap B \neq \emptyset$.
Thus $\mathcal{F}(G)$ and $\mathcal{F}(H)$ are cross-intersecting.\\ \\
$(ii) \Rightarrow (iii)$\\
Suppose that $\mathcal{F}(G)$ and $\mathcal{F}(H)$ are cross-intersecting. For simplicity, for each $l \in [n]$, we denote $x_{l}=\mu_{G}(l), y_{l}=\mu_{H}(l), z_{l}=x_{l}+y_{l}$. We have the following simple observations. We have $x_{l+1}=x_{l}+1$ if $l+1 \in G$ and $x_{l+1}=x_{l}$ if $l+1 \notin G$. Similarly, we also have $y_{l+1}=y_{l}+1$ if $l+1 \in H$ and $y_{l+1}=y_{i}$ if $l+1 \notin H$. Therefore, $z_{l+1}-z_{l} \in \{0,1,2\}$.\\
Assume, for a contradiction, that $\mu_{G}(l)+\mu_{H}(l) \leq l$ for all $1 \leq l \leq n$. Thus, $z_{l} \leq l$, for all $1 \leq l \leq n$. By induction, for each $l$, we will construct the sets $G_{l}$ and $H_{l}$ satisfying simultaneously the following conditions: $G_{l} \leq G \cap [l], H_{l} \leq H \cap [l], G_{l} \cap H_{l}= \emptyset$ and $G_{l} \cup H_{l}=[z_{l}]$.\\ 

Consider $l=1$. We have $z_{1} \leq 1$. So $z_{1}=0$ or $z_{1}=1$. If $z_{1}=0$ then we take $G_{1}=H_{1}= \emptyset$ and we see that the above conditions are trivially satisfied. We consider the case where $z_{1}=x_{1}+y_{1}=1$. We may assume that $x_{1}=1$ and $y_{1}=0$. This means that $1 \in G$ and $1 \notin H$. We take $G_{1}=\{1\} \leq G \cap [1]$ and $H_{1}= \emptyset \leq H \cap [1]$. The above conditions are also satisfied.\\
Assume that we built two sets $G_{l}$ and $H_{l}$ that satisfy the conditions $G_{l} \leq G \cap [l], H_{l} \leq H \cap [l], G_{l} \cap H_{l}= \emptyset$ and $G_{l}\cup H_{l}=[z_{l}]$. We will build the sets $G_{l+1}$ and $H_{l+1}$. There are three possible cases. \\
\textbf{ The case of $z_{l+1}=z_{l}$}\\
We have $l+1 \notin G\cup H$. We set $G_{l+1}=G_{l}$ and $H_{l+1}=H_{l}$. We have, by the induction hypothesis, $G_{l+1}=G_{l} \leq G \cap [l]= G \cap [l+1]$ and $H_{l+1}=H_{l} \leq H \cap [l]= H \cap [l+1]$. Furthermore, $G_{l+1} \cap H_{l+1}= G_{l} \cap H_{l}= \emptyset$ and $G_{l+1} \cup H_{l+1}=G_{l} \cup H_{l}=[z_{l}]=[z_{l+1}]$. \\
\textbf{ The case of $z_{l+1}=z_{l}+1$}\\
We have $x_{l+1}=x_{l}+1, y_{l+1}=y_{l}$ or $y_{l+1}=y_{l}+1, x_{l+1}=x_{l}$. Without losing generality, we can see that $x_{l+1}=x_{l}+1, y_{l+1}=y_{l}$. Thus, $l+1 \in G$ and $l+1 \notin H$. We set $G_{l+1}=G_{l} \cup \{z_{l+1}\}$ and $H_{l+1}=H_{l}$. According to the induction hypothesis, $G_{l} \leq G \cap [l]$ and $H_{l} \leq H \cap [l]$. Note that $z_{l+1} \leq l+1$. It implies $G_{l+1}=G_{l} \cup \{z_{l+1}\} \leq G\cap [l+1]$ and $H_{l+1}=H_{l} \leq H\cap [l]=H \cap [l+1]$.   Also, by the induction hypothesis, we have $G_{l+1} \cap H _{l+1}= (G_{l} \cup \{z_{l+1}\}) \cap H_{l}=(G_{l} \cap H_{l})\cup (H_{l} \cap \{z_{l+1}\}=\emptyset$ since $z_{l+1} \notin H_{l}$. We also have $G_{l+1} \cup H_{l+1}=(G_{l} \cup H_{l}) \cup \{z_{l+1}\}=[z_{l}] \cup \{z_{l+1}\}=[z_{l+1}]$.\\ 
\textbf{ The case of $z_{l+1}=z_{l}+2$}\\
In this case, we have $x_{l+1}=x_{l}+1$ and $ y_{l+1}=y_{l}+1$. So $l+1 \in G \cap H$. We set $G_{l+1} =G_{l} \cup \{z_{l}+1\}$ and $H_{l+1} =H_{l} \cup \{z_{l}+2\}$. We check all the conditions. First, $G_{l+1} \cup H_{l+1} = (G_{l} \cup \{z_{l+1}-1\}) \cup (H_{l} \cup \{z_{l+1}\})= (G_{l} \cup H_{l}) \cup \{z_{l+1}-1\}  \cup \{z_{l+1}\}=[z_{l}]\cup \{z_{i+1}-1\} \cup\{z_{l+1}\}=[z_{l+1}]$. Next, since $G_{i} \cap H_{l}= \emptyset$, we have $G_{l+1} \cap H_{l+1}= (G_{l} \cup \{z_{l}+1\}) \cap (H_{l} \cup \{z_{l}+2\})= \emptyset$. Finally, since $l+1 \in G \cap H$ and $G_{l} \leq G \cap [l]$, we have $G_{l+1}=G_{l} \cup \{z_{l}+1\} \leq (G \cap [l]) \cup \{l+1\}=G \cap [l+1]$. Similarly, we also have $H_{l+1} \leq H \cap [l+1]$.\\ 
Now, we let $m = \text{ max } \{g_{1},\ldots,g_{r}, h_{1},\ldots,h_{s}\}$. Note that $G \cap [m] =G, H \cap [m]=H$ and $m \leq |G|+|H|$. Applying the above construction for $l=m$, we get two sets $G_{m}, H_{m}$ satisfying the conditions: $G_{m} \leq G, H_{m} \leq H, G_{m} \cup H_{m}=[m]$ and $G_{m} \cap H_{m}=\emptyset$. We take two sets $K,J \subseteq[m+1,n]$ such that $|K|=k-|G|, |J|=k-|H|$ and $K \cap J=\emptyset$ (it is possible to choose two sets $K$ and $J$ because $k-|G|+k-|H| \leq n-m )$. We define $A=G_{m} \cup K$ and $B=H_{m} \cup J$. We have $|A|=|B|=k$, $A \preceq G$ and $B \preceq H$. Thus, $A \in \mathcal{F}(G)$ and $B \in \mathcal{F}(G)$. Since $\mathcal{F}(G)$ and $\mathcal{F}(H)$ are cross-intersecting, we have $A \cap B \neq \emptyset$.
However, this leads to a contradiction, since by the construction of the sets $A$ and $B$, we clearly have $A \cap B =\emptyset$. Thus, there exists $l \in [n]$ such that $\mu_{G}(l)+\mu_{H}(l) > l$ as desired.\\
$(iii) \Rightarrow (i)$\\
Assume that there exists $l \in [n]$ such that $\mu_{G}(l) +\mu_{H}(l) > l$. We need to prove that $G$ and $H$ are strongly intersecting. Take any $G' \leq G$ and $H' \leq H$. By Proposition 2.1, we have $\mu_{G'}(l) \geq \mu_{G}(l)$ and $\mu_{H'}(l) \geq \mu_{H}(l)$. Thus, $\mu_{G'}(l) +\mu_{H'}(l) \geq \mu_{G}(l) +\mu_{H}(l) > l$. This implies that $(G' \cap [l])\cap (H' \cap [l]) \neq \emptyset$ and therefore, $G' \cap H' \neq \emptyset$. Thus, $G \sim_{si} H$. 
\end{proof}

\end{document}
